\documentclass[11pt,a4paper]{article}

\usepackage[utf8]{inputenc}
\usepackage[T1]{fontenc}
\usepackage{geometry}
\usepackage{listings}
\usepackage{xcolor}
\usepackage{amsmath}
\usepackage{amssymb}
\usepackage{hyperref}
\usepackage{enumitem}
\usepackage{titlesec}

\geometry{margin=2.5cm}

% ---------- Estilo de codigo ----------
\definecolor{codebg}{rgb}{0.97,0.97,0.97}
\definecolor{codegray}{rgb}{0.5,0.5,0.5}
\definecolor{codeblue}{rgb}{0.0,0.0,0.6}
\definecolor{codegreen}{rgb}{0.0,0.45,0.0}

\lstdefinestyle{codigo}{
    backgroundcolor=\color{codebg},
    basicstyle=\ttfamily\small,
    commentstyle=\color{codegreen},
    keywordstyle=\color{codeblue}\bfseries,
    numberstyle=\tiny\color{codegray},
    numbers=left,
    numbersep=6pt,
    breaklines=true,
    breakatwhitespace=false,
    frame=single,
    rulecolor=\color{codegray},
    tabsize=2,
    showstringspaces=false,
    xleftmargin=1.2em,
    framexleftmargin=1em
}
\lstset{style=codigo}

\titleformat{\subsubsection}{\normalsize\bfseries}{\thesubsubsection}{1em}{}

\hypersetup{
    colorlinks=true,
    linkcolor=blue,
    urlcolor=blue,
    citecolor=blue
}

\title{\textbf{Documentaci\'on T\'ecnica del Compilador C-TDS}\\
\large Analizador L\'exico, Sint\'actico y Representaci\'on del AST}
\author{Taller de Dise\~no de Software (Cod. 3306) --- TDS26}
\date{Septiembre de 2026}

\begin{document}

\maketitle

\begin{abstract}
\noindent
Este documento describe el estado actual del compilador de C-TDS: el analizador
l\'exico (\texttt{lexer.l}), el analizador sint\'actico (\texttt{bison.y}) y la
representaci\'on del \'Arbol de Sintaxis Abstracta (\texttt{ast.h} / \texttt{ast.c}).
El objetivo no es solo describir \emph{qu\'e} hace cada archivo, sino \emph{por qu\'e}
est\'a hecho as\'i: qu\'e parte de la letra del proyecto o de la gram\'atica del
lenguaje motiv\'o cada decisi\'on, y qu\'e problemas concretos aparecieron (y c\'omo
se resolvieron) al implementarlo. La \'ultima secci\'on est\'a pensada espec\'ificamente
para que el resto del equipo pueda continuar agregando reglas a \texttt{bison.y}
siguiendo el mismo patr\'on que las que ya est\'an resueltas.
\end{abstract}

\tableofcontents

% ======================================================================
\section{Introducci\'on}
% ======================================================================

El compilador se construye de forma incremental: primero el an\'alisis l\'exico y sint\'actico, luego el sem\'antico,
y reci\'en despu\'es la generaci\'on de c\'odigo intermedio y objeto. Este
documento cubre las dos primeras etapas, m\'as la estructura de datos
(el AST) que las conecta con las etapas siguientes.

El lenguaje fuente es C-TDS, especificado en el documento
\emph{Descripci\'on del Lenguaje: C-TDS}. Todas las decisiones de dise\~no
que no est\'an fijadas expl\'icitamente por ese documento se marcan como
tales a lo largo de esta gu\'ia, junto con su justificaci\'on.

% ======================================================================
\section{Analizador L\'exico (\texttt{lexer.l})}
% ======================================================================

El analizador l\'exico convierte el texto fuente en una secuencia de
\emph{tokens}. Implementa exactamente el l\'exico definido en
\emph{Descripci\'on del Lenguaje: C-TDS}, secci\'on ``Consideraciones del
L\'exico''.

\subsection{Palabras reservadas}

La letra define expl\'icitamente el conjunto cerrado de palabras
reservadas:

\begin{quote}
\emph{``Las palabras reservadas son: boolean else false if int return
true void while float''}
\end{quote}

Cada una se reconoce con una regla propia que retorna el token
correspondiente (\texttt{BOOLEAN}, \texttt{ELSE}, \texttt{IF}, etc.).
Un punto importante: \textbf{\texttt{main} no es palabra reservada}.
Aunque todo programa v\'alido debe tener un m\'etodo \texttt{main}, esa
es una regla \emph{sem\'antica} (regla \#3 de la secci\'on ``Reglas
Sem\'anticas''): \emph{``Todo programa contiene la definici\'on de un
m\'etodo llamado main''}. No es una regla l\'exica ni sint\'actica, as\'i que
\texttt{main} matchea como un \texttt{ID} com\'un, y su existencia se
verificar\'a m\'as adelante contra la tabla de s\'imbolos, en la etapa de
an\'alisis sem\'antico.

\subsection{Identificadores}

La gram\'atica del lenguaje define:
\[
\langle id \rangle \rightarrow \langle alpha \rangle\, \langle alpha\_num \rangle^{*}
\qquad
\langle alpha\_num \rangle \rightarrow \langle alpha \rangle \mid \langle digit \rangle
\]

Es decir, un identificador es una letra seguida de cero o m\'as
letras/d\'igitos --- \textbf{sin gui\'on bajo}. La regla actual en
\texttt{lexer.l} es:

\begin{lstlisting}[language=C]
ALPHA_NUM   {ALPHA}|{DIGIT}|"_"
\end{lstlisting}

Este \texttt{"\_"} es una \textbf{extensi\'on} sobre la gram\'atica
estricta. La letra permite expl\'icitamente extensiones documentadas
(\emph{Descripci\'on del Proyecto}, secci\'on ``Documentaci\'on'', punto~2:
\emph{``Una lista de decisiones de dise\~no, asumpciones, aclaraciones o
extensiones realizadas en la etapa''}), as\'i que esto es v\'alido siempre
que quede documentado como tal --- que es lo que se hace ac\'a. Si el
equipo decide no permitir guiones bajos, basta con sacar el
\texttt{|"\_"} de esa l\'inea.

\subsection{N\'umeros: enteros y reales}

La gram\'atica separa dos literales num\'ericos:
\[
\langle int\_literal \rangle \rightarrow \langle digit \rangle\, \langle digit \rangle^{*}
\qquad
\langle float\_literal \rangle \rightarrow \langle digit \rangle\, \langle digit \rangle^{*}\, \text{.}\, \langle digit \rangle\, \langle digit \rangle^{*}
\]

Notar que \emph{ambos} lados del punto exigen al menos un d\'igito. Esto
se traduce directamente a dos definiciones separadas:

\begin{lstlisting}[language=C]
NUMBER      {DIGIT}+
FLOAT       {DIGIT}+"."{DIGIT}+
\end{lstlisting}

con la regla de \texttt{FLOAT} ubicada \textbf{antes} que la de
\texttt{NUMBER} en el archivo. Esto importa: flex resuelve ambig\"uedades
por longitud de coincidencia (\emph{maximal munch}), pero ante un
empate usa la regla que aparece primero en el archivo. Poner
\texttt{FLOAT} primero evita cualquier duda sobre qui\'en gana ese
posible empate.

Como consecuencia directa de exigir d\'igitos a ambos lados del punto,
entradas como \texttt{5.} o \texttt{.5} \textbf{no} son floats v\'alidos:
se tokenizan como \texttt{NUMBER} m\'as un s\'imbolo \texttt{.} inv\'alido
suelto. Esto se verific\'o expl\'icitamente con casos de prueba y es el
comportamiento esperado seg\'un la gram\'atica, no un error.

\subsection{Comentarios}

La letra define dos formas de comentario:

\begin{quote}
\emph{``Dos tipos de comentarios son permitidos, aquellos que comienzan
con // y terminan al final de la l\'inea (\'unicamente pueden ser de una
l\'inea) y aquellos que est\'an delimitados por /* y */ (pueden tener
varias l\'ineas de extensi\'on).''}
\end{quote}

Ambos se descartan sin generar token, en l\'inea con: \emph{``S\'imbolos
que no son tokens son descartados, por ejemplo, espacios, tabulaciones,
nuevas l\'ineas y comentarios''}.

\subsection{N\'umero de l\'inea en los mensajes}

Se activ\'o \texttt{\%option yylineno} y se agreg\'o \texttt{[Linea \%d]}
a todos los mensajes de error. La letra no lo pide con esas palabras
exactas, pero s\'i exige, en la secci\'on ``Consideraciones del L\'exico'':

\begin{quote}
\emph{``Aquellos s\'imbolos que no son permitidos en el lenguaje deben ser
reportados, por ejemplo, \$, \#.''}
\end{quote}

Reportar un s\'imbolo inv\'alido sin decir en qu\'e l\'inea aparece es de
poca utilidad pr\'actica en un archivo real. Adem\'as, la siguiente etapa
(el analizador sint\'actico) necesita la l\'inea para sus propios mensajes
de error --- \texttt{yyerror} en \texttt{bison.y} ya usa
\texttt{yylineno}, que \texttt{lexer.l} mantiene actualizado. Es una
decisi\'on de dise\~no, documentada como tal.

\subsection{Errores l\'exicos}

La regla de error (\texttt{.} como \'ultima alternativa) reporta
\texttt{ERROR LEXICO: simbolo no permitido} en vez del gen\'erico
\texttt{Token inv\'alido} usado en versiones anteriores, para que el
mensaje quede alineado con la terminolog\'ia de la letra citada arriba.

% ======================================================================
\section{Analizador Sint\'actico (\texttt{bison.y})}
% ======================================================================

El analizador sint\'actico verifica que la secuencia de tokens producida
por el lexer respete la gram\'atica de C-TDS, y construye el AST
correspondiente a medida que reconoce cada producci\'on.

\subsection{Estado actual: alcance deliberadamente reducido}

\textbf{Este archivo no implementa toda la gram\'atica del lenguaje.}
Implementa \'unicamente 4 producciones, elegidas para servir de patr\'on
de referencia:

\begin{enumerate}
    \item \texttt{Type} --- nodo hoja, sin hijos.
    \item \texttt{VarDecl} (con \texttt{IdList}) --- nodo con un hijo
        fijo m\'as una cantidad variable de hijos.
    \item \texttt{Statement} (caso asignaci\'on) --- nodo binario con
        \texttt{left}/\texttt{right}.
    \item \texttt{Expr} (caso \texttt{NUMBER}) --- nodo hoja con un
        valor literal.
\end{enumerate}

El resto de la gram\'atica (\texttt{Program} real, \texttt{MethodDecl},
\texttt{Block}, \texttt{if}/\texttt{while}, llamadas a m\'etodo, y el
resto de los operadores de \texttt{Expr}) \textbf{no est\'a declarado
todav\'ia}. 

El s\'imbolo inicial actual es un \texttt{Program} de prueba m\'inimo
(\texttt{VarDecl Statement}), \'util \'unicamente para poder compilar y
probar las 4 reglas ya escritas. \textbf{No es la gram\'atica real del
lenguaje} y se va a reemplazar cuando se implemente el \texttt{Program}
verdadero (secci\'on~\ref{sec:program-real}).

\subsection{Por qu\'e \texttt{main} no es un token}

Se sac\'o deliberadamente un token \texttt{MAIN} que exist\'ia en una
versi\'on anterior del archivo. Como se explic\'o en la secci\'on del
lexer, la existencia de un m\'etodo \texttt{main} es una regla
sem\'antica, no sint\'actica: \texttt{main} debe llegar al parser como un
\texttt{ID} normal, igual que cualquier otro identificador.

\subsection{Un problema real encontrado: ambig\"uedad \texttt{Type ID}}
\label{sec:program-real}

Al dise\~nar c\'omo ser\'ia el \texttt{Program} real, apareci\'o un problema
concreto que vale la pena documentar en detalle, porque quien complete
esta regla se va a topar con \'el si transcribe la gram\'atica de forma
literal.

La gram\'atica de \emph{Descripci\'on del Lenguaje: C-TDS} define:
\[
\langle program \rangle \rightarrow \langle var\_decl \rangle^{*}\; \langle method\_decl \rangle^{*}
\]

Una traducci\'on directa a Bison usar\'ia dos no-terminales separados,
uno para la lista de \texttt{var\_decl} y otro para la lista de
\texttt{method\_decl}. El problema es que \textbf{ambas producciones
empiezan igual}:

\[
\langle var\_decl \rangle \rightarrow \langle type \rangle\; \langle id \rangle^{+,}\; \text{;}
\qquad
\langle method\_decl \rangle \rightarrow \langle type \rangle\; \langle id \rangle\; \text{(} \ldots \text{)}\; \langle block \rangle
\]

Un parser LALR(1) --- como el que genera Bison --- decide qu\'e regla
aplicar mirando \'unicamente \emph{un} s\'imbolo por delante
(\emph{lookahead}). Al ver \texttt{Type ID}, no hay forma de saber
todav\'ia si lo que sigue es \texttt{;}, \texttt{,} (entonces es una
variable) o \texttt{(} (entonces es un m\'etodo) sin seguir leyendo. Al
declarar \texttt{VarDecl} y \texttt{MethodDecl} como listas separadas,
Bison reporta un conflicto \emph{shift/reduce} y lo resuelve siempre a
favor de la variable --- con el resultado de que \textbf{ning\'un
m\'etodo con tipo de retorno distinto de \texttt{void} llega a
parsear jam\'as}. Este bug se detect\'o probando expl\'icitamente contra
el ejemplo de programa que trae la propia letra del lenguaje.

\paragraph{Soluci\'on.} Unificar ambas producciones bajo un \'unico
no-terminal (por ejemplo \texttt{Declaration}, con
\texttt{DeclList} como la lista que acumula ambos casos indistintamente)
y decidir reci\'en \emph{despu\'es} de leer el identificador si contin\'ua
como declaraci\'on de variable o de m\'etodo. Esto no cambia el lenguaje
reconocido --- sigue aceptando exactamente los mismos programas --- pero
elimina la ambig\"uedad estructural. Es la soluci\'on est\'andar para
prefijos compartidos en gram\'aticas LL/LALR.

\textbf{Qui\'en implemente \texttt{Program}/\texttt{MethodDecl}/\texttt{Block}
debe tener esto presente} y no separar \texttt{VarDecl} de
\texttt{MethodDecl} como dos listas independientes al nivel de
programa. (Dentro de un \texttt{Block} s\'i es seguro usar solo
\texttt{VarDecl}, porque la gram\'atica no permite declarar m\'etodos
dentro de un bloque.)

\subsection{Precedencia de operadores}

La tabla de precedencia de \texttt{bison.y} reproduce exactamente la de
\emph{Descripci\'on del Lenguaje: C-TDS}, secci\'on ``Expresiones''
(de menor a mayor precedencia):

\begin{center}
\begin{tabular}{ll}
\hline
Operadores & Asociatividad \\
\hline
\texttt{||} & izquierda \\
\texttt{\&\&} & izquierda \\
\texttt{==} & izquierda \\
\texttt{<}\ \texttt{>} & izquierda \\
\texttt{+}\ \texttt{-} & izquierda \\
\texttt{*}\ \texttt{/}\ \texttt{\%} & izquierda \\
\texttt{!}\ (negaci\'on),\ \texttt{-}\ (unario) & derecha \\
\hline
\end{tabular}
\end{center}

El \texttt{\%prec UMINUS} se usa (cuando se implemente la regla de menos
unario en \texttt{Expr}) para distinguir la resta binaria del menos
unario, que comparten el mismo s\'imbolo \texttt{-} pero tienen
precedencias distintas.

% ======================================================================
\section{Representaci\'on del AST (\texttt{ast.h} y \texttt{ast.c})}
% ======================================================================

\subsection{Origen de la estructura: \texttt{redesign\_compilator.pdf}}

El dise\~no de \texttt{NodeAST} y \texttt{Symbol} sigue el documento de
dise\~no del equipo:

\begin{lstlisting}[language=C]
typedef struct Symbol {
    char *id;
    char *value;
} Symbol;

typedef struct NodeAST {
    Symbol *symbol;
    struct NodeAST *left;
    struct NodeAST *right;
    DataType type;
    NodeType nodeType;
} NodeAST;
\end{lstlisting}

Dos decisiones de ese documento vale la pena resaltar porque no son
obvias a primera vista:

\begin{itemize}
    \item \textbf{El tipo no vive en \texttt{Symbol}, vive en
        \texttt{NodeAST.type}.} Esto evita tener dos fuentes de verdad
        distintas para el tipo de una expresi\'on --- si estuviera en
        ambos lugares, podr\'ian desincronizarse.
    \item \textbf{\texttt{Symbol} no es exclusivamente para variables.}
        Tambi\'en se usa para representar operadores o resultados
        temporales durante la resoluci\'on del AST (por ejemplo, un
        \texttt{BINOP\_NODE} para \texttt{a + b} puede tener
        \texttt{symbol = \{id: "+", value: "t1"\}}, donde \texttt{t1}
        es un temporal). Esto es relevante para cuando se implemente
        el generador de c\'odigo intermedio: los temporales de tres
        direcciones ya tienen un lugar natural donde vivir.
\end{itemize}

\subsection{Qu\'e se agreg\'o sobre ese documento, y por qu\'e}

\texttt{redesign\_compilator.pdf} ejemplific\'o la resoluci\'on de
identificadores con un subconjunto reducido de casos (asignaci\'on,
operaci\'on binaria, constante, retorno, definici\'on). La gram\'atica
completa del lenguaje necesita m\'as que eso. Concretamente, sobre
\texttt{NodeAST} se agreg\'o:

\begin{lstlisting}[language=C]
struct NodeAST **children;
int childCount;
\end{lstlisting}

porque varias producciones de la gram\'atica tienen una cantidad
\textbf{variable} de hijos, no fija en dos (\texttt{left}/\texttt{right}
alcanza para nodos binarios, pero no para un programa con $N$
declaraciones, un m\'etodo con $N$ par\'ametros, o una llamada con $N$
argumentos). El campo \texttt{TYPE\_FLOAT} en \texttt{DataType} tambi\'en
se agreg\'o porque el documento original solo contemplaba
\texttt{INT}/\texttt{BOOL}/\texttt{VOID}, mientras que la letra del
lenguaje define expl\'icitamente: \emph{``Los tipos b\'asicos en C-TDS son
int, float y boolean''}.

\textbf{Estado actual (versi\'on m\'inima):} el \texttt{enum NodeType}
en \texttt{ast.h} solo declara los 5 valores que las 4 reglas ya
implementadas en \texttt{bison.y} necesitan:
\texttt{DEFINITION\_NODE}, \texttt{VAR\_DECL\_NODE}, \texttt{ID\_NODE},
\texttt{ASSIGNMENT\_NODE} y \texttt{CONSTANT\_NODE}. A medida que se
implementen m\'as reglas de la gram\'atica, va a hacer falta agregar
los \texttt{NodeType} correspondientes, por ejemplo:
\begin{itemize}[nosep]
    \item \texttt{METHOD\_DECL\_NODE}, \texttt{BLOCK\_NODE}
    \item \texttt{IF\_NODE}, \texttt{WHILE\_NODE}
    \item \texttt{METHOD\_CALL\_NODE}, \texttt{UNARYOP\_NODE}, \texttt{PARAM\_NODE}
\end{itemize}
Se decide en el momento de implementar cada regla, no antes.

\subsection{\texttt{NodeList}: por qu\'e existe y por qu\'e no es parte del AST final}

Varias producciones de la gram\'atica son listas (\emph{cero o m\'as},
\emph{uno o m\'as separados por coma}). Bison construye el resultado de
una regla \emph{de abajo hacia arriba}: cuando se reduce
\texttt{IdList: IdList ',' ID}, todav\'ia no existe el nodo padre
(\texttt{VarDecl}) que va a contener esa lista como hijos. Hace falta
alg\'un lugar transitorio donde acumular esos nodos hermanos mientras se
van reconociendo, uno por uno, de izquierda a derecha.

\texttt{NodeList} cumple exactamente ese rol: es una lista enlazada
simple, usada \'unicamente durante la construcci\'on (dentro de las
acciones de \texttt{bison.y}). Una vez que el nodo padre existe,
\texttt{attachChildren} vuelca esa lista al arreglo \texttt{children[]}
del padre y fija \texttt{childCount}. Ning\'un \texttt{NodeAST} del
\'arbol final contiene un \texttt{NodeList}: es puramente un andamiaje de
construcci\'on, no una estructura de datos del AST en s\'i.

\subsection{Las cuatro funciones, una por una}

\subsubsection{\texttt{newSymbol}}

\begin{lstlisting}[language=C]
Symbol *newSymbol(const char *id, const char *value);
\end{lstlisting}

Crea un \texttt{Symbol} reservado din\'amicamente. El detalle de
implementaci\'on m\'as importante: \textbf{debe copiar} las cadenas
recibidas (t\'ipicamente con \texttt{strdup}), no quedarse con el
puntero original. La raz\'on es que \texttt{id} normalmente viene de
\texttt{yytext} (el lexema que reconoci\'o flex), y flex \textbf{reutiliza
el mismo buffer} en cada token siguiente. Si \texttt{Symbol} guardara
el puntero crudo en vez de una copia, el contenido cambiar\'ia por
debajo sin aviso apenas el lexer procese el pr\'oximo token, corrompiendo
silenciosamente todos los \texttt{Symbol} creados hasta ese momento.

\subsubsection{\texttt{newNode}}

\begin{lstlisting}[language=C]
NodeAST *newNode(NodeType nodeType, Symbol *symbol,
                  NodeAST *left, NodeAST *right);
\end{lstlisting}

Crea un \texttt{NodeAST}, dejando \texttt{children == NULL} y
\texttt{childCount == 0} (se completan aparte, ver
\texttt{attachChildren}, solo en los nodos que los necesitan). El
campo \texttt{type} se espera en un valor por defecto: el tipo
\emph{real} de la mayor\'ia de los nodos no se conoce en el momento
sint\'actico --- reci\'en se resuelve en el an\'alisis sem\'antico, cuando
se puede consultar la tabla de s\'imbolos para saber, por ejemplo, con
qu\'e tipo fue declarada una variable. Los pocos casos donde el tipo
\emph{s\'i} se conoce de entrada (una constante entera es siempre
\texttt{TYPE\_INT}, un \texttt{Type: int} siempre da \texttt{TYPE\_INT})
lo fijan expl\'icitamente despu\'es de llamar a \texttt{newNode}, como se
ve en las reglas \texttt{Type} y \texttt{Expr: NUMBER} de
\texttt{bison.y}.

\subsubsection{\texttt{newNodeList}}

\begin{lstlisting}[language=C]
NodeList *newNodeList(NodeAST *node, NodeList *next);
\end{lstlisting}

Crea (o extiende) una celda de la lista transitoria descripta arriba.
El patr\'on de uso en \texttt{bison.y}, visible en la regla
\texttt{IdList}, es: crear una lista de un elemento en el primer caso
base, y en el caso recursivo recorrer hasta el final de la lista
existente para agregar el nuevo elemento, preservando el orden textual
en que aparecen en el c\'odigo fuente.

\subsubsection{\texttt{attachChildren}}

\begin{lstlisting}[language=C]
void attachChildren(NodeAST *parent, NodeList *list);
\end{lstlisting}

Centraliza el patr\'on ``contar los elementos de la lista, reservar el
arreglo del tama\~no exacto, copiar cada nodo en orden'' que, sin esta
funci\'on, se repetir\'ia casi id\'entico en cada regla de la gram\'atica
que tenga una cantidad variable de hijos. Se us\'o en \texttt{VarDecl} y
se va a volver a usar en \texttt{MethodDecl} (para los par\'ametros),
\texttt{MethodCall} (para los argumentos) y \texttt{Program} (para la
lista de declaraciones), entre otras.

\subsection{Por qu\'e las cuatro funciones tienen el cuerpo vac\'io}

Esto es deliberado, no un olvido. Se decidi\'o definir primero la
\textbf{interfaz completa} (qu\'e recibe cada funci\'on, qu\'e devuelve, qui\'en
es due\~no de la memoria en cada caso --- todo documentado en los
comentarios de \texttt{ast.h}) y reci\'en despu\'es completar la
implementaci\'on real en \texttt{ast.c}. Esto permite que \texttt{bison.y}
ya compile contra una interfaz estable mientras la implementaci\'on se
sigue discutiendo o repartiendo entre el equipo, sin bloquear el
trabajo de nadie.

\end{document}
